\documentclass[11pt]{article}
\usepackage[T1]{fontenc}
\usepackage{lmodern}
\usepackage[margin=1in]{geometry}
\usepackage{amsmath,amssymb,amsthm,mathtools}
\usepackage{microtype}
\usepackage[colorlinks=true,linkcolor=blue,citecolor=blue,urlcolor=blue]{hyperref}
\hypersetup{pdfauthor={OpenAI},pdftitle={Weak mixing of triangular billiards with an irrational angle}}
\newtheorem{theorem}{Theorem}[section]
\newtheorem{lemma}[theorem]{Lemma}
\newtheorem{proposition}[theorem]{Proposition}
\newtheorem{corollary}[theorem]{Corollary}
\theoremstyle{remark}
\newtheorem{remark}[theorem]{Remark}
\newcommand{\R}{\mathbb R}
\newcommand{\C}{\mathbb C}
\newcommand{\Z}{\mathbb Z}
\newcommand{\eps}{\varepsilon}
\newcommand{\pa}{\partial}
\newcommand{\pb}{\bar\partial}
\newcommand{\ind}{\mathbf 1}
\newcommand{\norm}[1]{\left\lVert#1\right\rVert}
\newcommand{\ip}[2]{\left\langle#1,#2\right\rangle}
\setlength{\emergencystretch}{2em}
\title{Weak mixing of triangular billiards\\with an irrational angle}
\author{OpenAI}
\date{October 5, 2026}
\begin{document}
\maketitle
\begin{abstract}
We prove that the unit-speed billiard flow in every nondegenerate Euclidean
triangle with at least one angle irrational relative to \(\pi\) is weakly
mixing for normalized Liouville measure. Equivalently, the product of the
flow with itself is ergodic.
\end{abstract}
\section{Introduction}

Let \(Q\subset\R^2\) be a nondegenerate Euclidean triangle. Its unit-speed
billiard flow \(\Phi_t\) moves in straight lines in the interior and reflects
specularly at the open sides. Trajectories meeting a vertex in either time
direction are discarded; they form a null set. The invariant probability
measure on position and direction is normalized Liouville measure
\[
 d\mu(x,\theta)=\frac{dA(x)\,d\theta}{2\pi\operatorname{Area}(Q)}.
\]
We say that the flow is \emph{weakly mixing} if its product with itself,
\((x,y)\mapsto(\Phi_t x,\Phi_t y)\), is ergodic for
\(\mu\otimes\mu\). Thus every measurable set invariant under all product
time maps has measure zero or one.

\begin{theorem}\label{thm:main}
The unit-speed billiard flow in every nondegenerate Euclidean triangle with
at least one angle irrational relative to \(\pi\) is weakly mixing for
normalized Liouville measure.
\end{theorem}

In particular, for triangles of area one with angles
\(\alpha,\beta,\pi-\alpha-\beta\), the conclusion holds for
Lebesgue-almost every point of
\[
 D=\{(\alpha,\beta)\in\R^2:\alpha>0,\ \beta>0,\ \alpha+\beta<\pi\}:
\]
the excluded pairs with both \(\alpha/\pi\) and \(\beta/\pi\) rational
form a countable set. The irrationality assumption is sharp for the full
phase space. In a rational triangle, side reflections generate a finite
group on directions. A nonconstant function on the circle invariant under
this finite group is then an invariant function for the billiard flow, so
the full flow is not ergodic.

\paragraph{Historical context.}
Polygonal billiards connect elementary reflection geometry with dynamics
on flat surfaces. The unfolding construction of Zemlyakov and Katok
\cite{ZK} replaces reflections by straight continuation on a surface.
For rational polygons, it produces a compact translation surface and
separates the full phase space into invariant directional surfaces.
Kerckhoff, Masur, and Smillie proved unique ergodicity in almost every
direction \cite{KMS}. They also obtained ergodicity of the full billiard
flow for a dense \(G_\delta\) family of polygons. Vorobets gave a
quantitative rational-approximation criterion and explicit ergodic
examples, including irrational right triangles \cite{Vorobets}.
These results establish ergodicity under genericity or sufficiently rapid
approximation assumptions rather than for each irrational triangle.

Chaika and Forni proved that the full billiard flow is weakly mixing for a
dense \(G_\delta\) set of polygons with any fixed number of sides at least
three \cite{CF}. Their argument uses the absence of common nontrivial
eigenvalues for directional flows in almost every pair of directions on a
translation surface. The Baire-category conclusion does not assert
Lebesgue-almost-everywhere weak mixing in the space of triangle angles.
For rational billiards, recent work of Arana-Herrera, Chaika, and Forni
characterizes weak mixing in almost every direction \cite{ACF}; this is
a statement on directional invariant surfaces, rather than on the full
Liouville phase space considered here.

Numerical investigations have long suggested weak mixing in irrational
right triangles and mixing in triangles with all angles irrational
\cite{ACG,CP}. Other computations have suggested nonergodicity in some
irrational right or symmetric triangles \cite{WCP,ZSBJ}. These conflicting
finite-time observations underscore the difficulty of identifying
measurable invariants from trajectories with slow statistical convergence.
Theorem~\ref{thm:main} treats every irrational triangle without an
approximation or symmetry restriction.

\paragraph{Analytic contribution.}
The regularity of an eigenfunction is the central issue. Glue two oppositely
oriented copies of \(Q\) along their sides. The resulting flat double is
regular across the open sides and has cone points at the vertices. On its
unit tangent bundle, let \(X\) denote differentiation along a geodesic and
\(Y\) differentiation in the perpendicular parallel direction. Locally
these vector fields commute. Forni and Moll developed angular Fourier and
Cauchy--Riemann methods in this setting \cite[Sections~3.2--3.5]{FM}.
Their rigidity results distinguish invariance under both horizontal
directions from geodesic invariance under an additional horizontal Sobolev
hypothesis \cite[Theorem~4.1 and Corollary~4.2]{FM}.
A measurable flow eigenfunction supplies only \(Xf=i\lambda f\) at the
outset; the missing transverse derivative cannot be assumed.

The companion manuscript \cite[Sections~2--6]{Companion} removes this
regularity obstacle for invariant functions, proving ergodicity of
irrational triangular billiards. We extend its coefficient-energy and
long-average method to every real spectral parameter \(\lambda\).
Summing the coefficient-energy identities over consecutive angular indices
makes the terms containing \(\lambda\) telescope to a single boundary term.
A separate pairing with the two boundary modes produced by differentiating
an angular Fourier tail gives the sharp bound
\[
 \norm{Xf}_2^2+\norm{Yf}_2^2\le\lambda^2\norm{f}_2^2.
\]
Since \(Xf=i\lambda f\), it follows that \(Yf=0\).
The use of a spectral parameter in Cauchy--Riemann energy estimates has a
related precedent in the twisted cohomological equations studied by Forni
\cite[Section~3]{Forni2022}. Here the argument starts with bounded measurable
eigenfunctions and proves the needed derivative bounds. We give a
self-contained proof of this analytic extension below.

\paragraph{Proof strategy.}
Expand a bounded eigenfunction on the double in angular Fourier modes.
Spatial smoothing away from the cone points and cutoffs near the tips
first give a uniform bound on the gradient of each coefficient. These
individual bounds do not
yet give a square-summable family of gradients. To obtain the missing
summability, take the sum of all modes above a fixed index. Its derivative
under \(X-i\lambda\) consists of two endpoint terms.

We test this equation using long time averages supported on orbit segments
that stay clear of the cone points. There are two approximation parameters.
At a fixed bound on the discarded measure, the smoothing limit gives a
bounded solution of the exact local eigenfunction equation and passes the
endpoint pairing to it.
The individual-coefficient estimate then supplies a new derivative bound
independent of the discarded measure. A second limit passes the pairing
to the original eigenfunction. Summing the resulting endpoint inequalities
proves \(Yf=0\). This order of limits is what permits the initial
eigenfunction to be merely measurable.

The remaining argument is geometric. On the triangle interior,
\(Xf=i\lambda f\) and \(Yf=0\) give a plane wave with a measurable
direction-dependent amplitude. Reflection in the two sides through an
irrational vertex makes this amplitude constant. Reflection in the third
side excludes every nonzero frequency. The standard spectral criterion
then gives weak mixing.

Section~\ref{sec:geometry} sets up the double and its Fourier coefficients.
Section~\ref{sec:energy} proves the coefficient and smoothing estimates.
Section~\ref{sec:rigidity} establishes transverse rigidity by the two
successive limits. Section~\ref{sec:conclusion} applies side reflections
and proves ergodicity of the product flow.

\section{The double and its angular coefficients}\label{sec:geometry}

We replace reflection by straight continuation on a flat surface. This
allows us to study the spatial derivatives of an eigenfunction without
imposing boundary conditions on the sides of the triangle. Throughout
Sections~\ref{sec:geometry}--\ref{sec:rigidity}, the triangle is fixed and
nondegenerate; its angles need not satisfy an irrationality assumption.

Glue two closed copies of the triangle \(Q\) along corresponding sides,
give the copies opposite orientations, and remove the three vertices.
The resulting connected oriented surface \(M\) has a flat metric. Across
an open side, a flat chart is obtained by reflecting one copy across
that side. Consequently, changes between oriented flat charts are
rotations followed by translations. The metric completion \(\overline M\)
is compact; its missing points form a set \(V\) of three cone points,
with cone angles twice the corresponding angles of \(Q\). This double
and the coefficient framework below are also used in
\cite[Section~2]{Companion}.

Let \(SM\) be the unit tangent bundle, let \(\nu\) be normalized
Liouville measure, and denote geodesic motion by \(\Psi_t\). Thus, in
a flat chart, \(\nu=d\bar A\,d\theta/(2\pi)\), where \(d\bar A\)
is probability area measure on \(M\). We first justify using this
motion as a measure-preserving flow defined modulo a null set.

A finite-time obstruction to a geodesic can occur only at a cone point.
Indeed, a unit-speed path is Cauchy as its time approaches a finite
endpoint, and its limit in \(\overline M\) permits continuation if it
is regular. There is also no unaccounted accumulation of side crossings
at a cone point. Lift the polar angle in a punctured cone to \(\R\)
and develop the lifted geodesic into the plane. Its image is a straight
line. If the path approaches the cone tip, this line ends at the
origin, so its final segment has constant polar angle and is radial.
Apart from trajectories along a seam, which
form a null set, that segment has no side crossings. On an earlier
compact regular segment crossings cannot accumulate: a flat chart at
a proposed accumulation point makes both the path and each local seam
straight. Thus every vertex hit has a finite crossing itinerary.
For each fixed finite itinerary, unfolding makes a vertex hit the
condition that a straight trajectory pass through a fixed point. For
almost every initial position this condition selects only finitely
many directions, so it has zero Liouville measure. Taking the countable
union over itineraries, vertices, and both time directions proves that
almost every geodesic exists for all real times.

Local geodesic motion is translation at fixed angle in flat charts and
preserves Liouville measure. Following a regular trajectory through
finitely many charts proves that every time map preserves \(\nu\),
with inverse \(\Psi_{-t}\) modulo null sets. Folding the two sheets to
\(Q\), and folding their tangent vectors by the corresponding local
isometries, defines a factor map to the billiard flow. It sends \(\nu\)
to \(\mu\), since the two sheets have equal area and the angular
measure is preserved by reflection. In particular, billiard
eigenfunctions lift to geodesic-flow eigenfunctions. Factor identities
are understood almost everywhere for each fixed time.

The composition operators \(u\mapsto u\circ\Psi_t\) form a strongly
continuous unitary group on \(L^2(SM,\nu)\). To see continuity at zero,
first take a continuous compactly supported function on the regular
bundle and use pointwise continuity of regular paths and dominated
convergence. Density then gives the assertion for every \(L^2\)
function. The same assertion for the billiard flow follows through
the isometric embedding given by the factor map.

In oriented flat coordinates \(x=(x_1,x_2)\), write
\(v=(\cos\theta,\sin\theta)\), and set
\[
 \pa=\tfrac12(\partial_{x_1}-i\partial_{x_2}),\qquad
 \pb=\tfrac12(\partial_{x_1}+i\partial_{x_2}).
\]
The horizontal derivatives in the forward and perpendicular directions
are
\begin{equation}\label{eq:XY}
 X=e^{i\theta}\pa+e^{-i\theta}\pb,\qquad
 Y=i e^{i\theta}\pa-i e^{-i\theta}\pb.
\end{equation}
These vector fields are globally defined on \(SM\), commute on the
regular bundle, and are formally skew adjoint for \(\nu\).
All distributional equations in this paper use compactly supported
smooth tests on \(M\) or \(SM\). No test passes through a cone point.

We use angular Fourier coefficients, as in the flat-flow framework of
Forni and Moll \cite[Sections~3.2--3.5]{FM}. For \(j\in\Z\), let
\(P_j\) be fiberwise Fourier projection, written in a chart as
\[
 P_j u=u_j(x)e^{ij\theta},\qquad
 u_j(x)=\int_0^{2\pi}u(x,\theta)e^{-ij\theta}\frac{d\theta}{2\pi}.
\]
Under an oriented change of coordinates \(z'=e^{i\phi}z+c\),
\(\theta'=\theta+\phi\), so
\[
 u'_j=e^{-ij\phi}u_j,\qquad
 (\pa u_j)'=e^{-i(j+1)\phi}\pa u_j,\qquad
 (\pb u_j)'=e^{-i(j-1)\phi}\pb u_j.
\]
Thus a coefficient is a section of a flat unitary line bundle, and its
two complex derivatives are coefficients of indices \(j+1\) and
\(j-1\), respectively. Taking Fourier coefficients in
\eqref{eq:XY} gives, also in distributions,
\begin{equation}\label{eq:modes}
 (Xu)_j=\pa u_{j-1}+\pb u_{j+1},\qquad
 (Yu)_j=i\pa u_{j-1}-i\pb u_{j+1}.
\end{equation}
The transition factors of the two terms in each sum agree.

We write \(\nabla\) for the flat spatial gradient, keeping the chart
angle fixed for a field on \(SM\). Coefficient gradients are weak
gradients in the preceding unitary trivializations. Their norms use
\(d\bar A\); bundle norms use \(\nu\). Whenever the horizontal
derivatives are square integrable, Parseval and the orthogonality of
the forward and perpendicular directions give
\begin{equation}\label{eq:gradient}
 \norm{\nabla u}_2^2
 =\norm{Xu}_2^2+\norm{Yu}_2^2
 =\sum_{j\in\Z}\norm{\nabla u_j}_2^2.
\end{equation}
For each coefficient, pointwise algebra also gives
\(\lvert\nabla u_j\rvert^2
=2(\lvert\pa u_j\rvert^2+\lvert\pb u_j\rvert^2)\).
An \(L^2\) gradient always means a finite norm on the regular surface;
it makes no assertion about traces or operator domains at a cone tip.
Inner products are complex and linear in the first entry.

Fix a real number \(\lambda\), and put \(A=X-i\lambda\), which is
again formally skew adjoint. If
\(f\circ\Psi_t=e^{i\lambda t}f\) in \(L^2\) for each \(t\), then
\(Af=0\) locally in distributions. Indeed, for a compactly supported
smooth test \(\varphi\), measure preservation gives
\[
 \ip{f}{\varphi\circ\Psi_{-t}}
 =e^{i\lambda t}\ip{f}{\varphi}.
\]
Differentiation at zero is justified on a compact regular neighborhood
of the test support and yields the claimed equation. Our analytic goal
is the following statement, whose proof occupies the next two sections.

\begin{proposition}[Transverse rigidity]\label{prop:rigidity}
Let \(M\) be the regular oriented double of a fixed nondegenerate
Euclidean triangle, and let \(\Psi_t\) be its geodesic flow on \(SM\)
with normalized Liouville measure. Suppose \(\lambda\in\R\) and
\(f\in L^\infty(SM)\) satisfy
\[
 f\circ\Psi_t=e^{i\lambda t}f
 \quad\text{in }L^2(SM)\quad\text{for every }t\in\R.
\]
Then \(f\) has square-integrable horizontal weak derivatives and
\[
 Yf=0,\qquad Xf=i\lambda f,\qquad
 \norm{\nabla f}_2^2=\lambda^2\norm{f}_2^2.
\]
\end{proposition}

The next section first controls the gradient of each angular
coefficient of a bounded local solution of \(Ah=0\).
Section~\ref{sec:rigidity} then obtains summability of these gradients
and proves the proposition.

\section{Localized energy and individual coefficient bounds}\label{sec:energy}

The first estimate uses the full spatial gradient only where the
eigenfunction equation fails. After smoothing and cutting off near
the vertices, this will confine the resulting error to small cone
annuli. The zero-frequency version of this argument appears in
\cite[Sections~3--4]{Companion}; the spectral correction below is
needed when \(\lambda\ne0\). The equal complex-derivative norms and
coefficient recurrences also occur in Forni and Moll's Sobolev
arguments \cite[Lemma~3.3, Corollary~4.2, and Lemma~5.1]{FM}.
Twisted complex-derivative energy identities provide a related
spectral precedent in Forni's work \cite[Section~3]{Forni2022}.

\begin{lemma}[Localized coefficient estimate]\label{lem:energy}
Let \(w\in L^2(SM)\) have compact spatial support in \(M\) and be
smooth in position with values in angular \(L^2\) in every flat chart.
Suppose \(Aw\) vanishes over \(M\setminus E\), where \(E\subset M\)
is measurable. Define
\(D_j=\norm{\pa w_j}_2^2\). Then, for every \(m\in\Z\),
\begin{equation}\label{eq:localized-energy}
 \tfrac12D_m+D_{m+1}
 \le \norm{Aw}_2\norm{\ind_E\nabla w}_2
       +\tfrac{\lambda^2}{2}\norm{w_{m+1}}_2^2,
 \qquad
 D_j=\tfrac14\norm{\nabla w_j}_2^2.
\end{equation}
\end{lemma}

\begin{proof}
Compact support in the regular surface permits integration by parts
between the coefficient bundles. The local formal adjoints are
\(\pa^*=-\pb\) and \(\pb^*=-\pa\); the derivatives commute in
flat charts. Hence
\begin{equation}\label{eq:complex-energies}
 \norm{\pa w_j}_2^2=\norm{\pb w_j}_2^2
                 =\tfrac14\norm{\nabla w_j}_2^2.
\end{equation}
These identities are intrinsic, since all transition factors are
constant and unitary. In particular, no integration takes place across
a cone point.

Put
\(C_j=\operatorname{Re}\ip{i\lambda w_{j+1}}{\pa w_j}\).
Equation~\eqref{eq:modes} gives
\begin{equation}\label{eq:spectral-telescope}
 D_j-D_{j+2}
 =\operatorname{Re}\ip{(Aw)_{j+1}}{\pa w_j-\pb w_{j+2}}
   +C_j-C_{j+1}.
\end{equation}
For the spectral term, the sign follows from
\[
 \operatorname{Re}\ip{i\lambda w_{j+1}}{\pb w_{j+2}}
 =\operatorname{Re}\ip{-i\lambda\pa w_{j+1}}{w_{j+2}}
 =C_{j+1}.
\]
The correction \(C_j-C_{j+1}\) couples consecutive indices, so a sum
over one parity class would not telescope. Sum
\eqref{eq:spectral-telescope} over \emph{every} integer
\(j=m,\ldots,N\). The left side is
\(D_m+D_{m+1}-D_{N+1}-D_{N+2}\), and the last two terms on the
right telescope to \(C_m-C_{N+1}\). All high endpoints tend to zero:
the coefficients of \(w\) and their spatial derivatives are square
summable, by the hypotheses on \(w\). The remaining sum of pairings
converges absolutely by Cauchy--Schwarz. Since \(Aw\) vanishes off
the part over \(E\), its absolute value is at most
\[
 \norm{Aw}_2
 \left(\int_E\sum_{j\ge m}
       |\pa w_j-\pb w_{j+2}|^2\,d\bar A\right)^{1/2}
 \le \norm{Aw}_2\norm{\ind_E\nabla w}_2.
\]
The last inequality uses the pointwise bound
\[
 \sum_{j\ge m}|\pa w_j-\pb w_{j+2}|^2
 \le 2\sum_{j\in\Z}(|\pa w_j|^2+|\pb w_j|^2)
 =\sum_{j\in\Z}|\nabla w_j|^2.
\]
Finally,
\[
 |C_m|\le |\lambda|\norm{w_{m+1}}_2\sqrt{D_m}
 \le \tfrac12D_m+\tfrac{\lambda^2}{2}\norm{w_{m+1}}_2^2,
\]
which proves \eqref{eq:localized-energy}.
\end{proof}

We now construct regularizations to which Lemma~\ref{lem:energy}
applies. Write \(r(x)=d(x,V)\) in \(\overline M\), also using \(r\)
for its pullback to \(SM\). Near each cone point, an embedded flat
disk about \(x\) has radius at least \(c r(x)\), for a fixed
\(c>0\): this follows by choosing uniformly sized disks on the compact
circle at radius one in the model cone and scaling. Compactness away
from the tips therefore gives a constant \(K>2\) such that, for all
sufficiently small \(\eps>0\), every point with
\(r>(K-2)\eps\) has an embedded flat disk of radius \(4\eps\).
For each fixed \(L>0\), the cone area formula also gives
\begin{equation}\label{eq:cone-area}
 \bar A\{r\le L\eps\}=O_L(\eps^2).
\end{equation}

Choose a nonnegative radial function
\(\rho\in C_c^\infty(\R^2)\), supported in the unit disk and with
Euclidean integral one, and set
\(\rho_\eps(y)=\eps^{-2}\rho(y/\eps)\). For \(r(x)>K\eps\),
define spatial convolution at fixed parallel direction by
\begin{equation}\label{eq:smoothing-definition}
 (S_\eps u)(x,\theta)
 =\int_{\R^2}\rho_\eps(y)u(x-y,\theta)\,dy
\end{equation}
in the embedded flat disk. Radiality makes this definition independent
of its oriented chart. Unless a derivative is being taken, extend
\(S_\eps u\) by zero to \(r\le K\eps\).

\begin{lemma}[Spatial regularization]\label{lem:smoothing}
For every integer \(\ell\ge0\), the preceding operators satisfy
\begin{equation}\label{eq:smoothing-bounds}
 \eps^\ell\norm{\nabla^\ell S_\eps u}_{L^2(r>K\eps)}
 \le C_\ell\norm{u}_2.
\end{equation}
Their derivatives here are taken before zero extension. The output
is smooth in position with values in angular \(L^2\), and
\(S_\eps u\to u\) in \(L^2(SM)\) for every fixed \(u\in L^2(SM)\).
For bounded input,
\[
 \norm{S_\eps u}_\infty\le\norm{u}_\infty,
 \qquad
 \eps^\ell\norm{\nabla^\ell S_\eps u}_{L^\infty(r>K\eps)}
 \le C_\ell\norm{u}_\infty.
\]
Convolution commutes locally with \(X,Y\), and \(A\), in
distributions. If \(L>K+4\) and \(g=0\) on \(r<L\eps\), then
\(S_\eps g=0\) on \(r<(L-1)\eps\); its zero extension has compact
spatial support in \(M\) and creates no boundary derivative.
In particular, when also \(Ag\in L^2\),
\begin{equation}\label{eq:smoothing-commutes}
 A(S_\eps g)=S_\eps(Ag)
\end{equation}
globally on the regular bundle, with the same zero-extension convention.
\end{lemma}

\begin{proof}
In a common flat neighborhood of nearby centers, this is ordinary
convolution with a kernel whose derivatives have the usual scaling.
Cauchy--Schwarz bounds the squared angular \(L^2\) norm of
\(\eps^\ell\nabla^\ell S_\eps u\) at a center by
\(C_\ell\eps^{-2}\) times the integral of
\(\norm{u(y,\cdot)}_{L^2_\theta}^2\) over its \(\eps\)-disk.
Parallel identification of tangent circles preserves angular measure.
A point in such a disk satisfies \(r>(K-1)\eps\), and all centers
which can contribute to it lie in a disk of area \(O(\eps^2)\).
Integrating the bound over centers proves
\eqref{eq:smoothing-bounds}. The same kernel estimates give the
\(L^\infty\) bounds; nonnegativity and mass one give the contraction
bound at order zero. Constant-coefficient differentiation in each
chart proves local commutation and spatial smoothness.

For a smooth compactly supported input on \(SM\), the usual
approximate-identity argument gives \(L^2\) convergence, including
the zero extension for sufficiently small \(\eps\). Such inputs
are dense in \(L^2(SM)\), so the uniform bound at order zero proves
convergence for every fixed input.

The function \(r\) is 1-Lipschitz. If \(r(x)<(L-1)\eps\), every
point sampled by its convolution disk has \(r<L\eps\), proving
the support claim. The convolution consequently vanishes in a whole
neighborhood of the boundary \(r=K\eps\), so zero extension produces
no boundary distribution. Local commutation gives
\eqref{eq:smoothing-commutes} there; on the omitted region both sides
vanish, since \(Ag\) has no support where \(g\) vanishes.
\end{proof}

\begin{lemma}[Individual coefficient regularity]\label{lem:bounded-modes}
Suppose \(h\in L^\infty(SM)\), \(\norm{h}_\infty\le H\), and
\((X-i\lambda)h=0\) locally in distributions on the regular bundle.
Every Fourier coefficient then has a global square-integrable flat
gradient, and
\begin{equation}\label{eq:bounded-modes}
 \norm{\nabla h_j}_2\le C_\lambda H\qquad(j\in\Z).
\end{equation}
The constant depends only on the fixed triangle and \(\lambda\).
\end{lemma}

\begin{proof}
Choose a smooth spatial cutoff \(\chi_\eps\), equal to zero on
\(r\le(K+1)\eps\), equal to one on \(r\ge(K+2)\eps\), and
satisfying \(0\le\chi_\eps\le1\) and
\(|\nabla\chi_\eps|\le C/\eps\). Its transition region
\(E_\eps\) consists of three disjoint cone annuli for small
\(\eps\). The field
\(w_\eps=\chi_\eps S_\eps h\), extended by zero near the tips,
has compact spatial support and the smoothness required by
Lemma~\ref{lem:energy}. Local commutation and \(Ah=0\) give
\[
 Aw_\eps=(X\chi_\eps)S_\eps h.
\]
This field is supported over \(E_\eps\), whose area is
\(O(\eps^2)\). On those annuli Lemma~\ref{lem:smoothing} gives
\[
 |Aw_\eps|+|\nabla w_\eps|\le CH/\eps.
\]
Consequently,
\(\norm{Aw_\eps}_2\le CH\) and
\(\norm{\ind_{E_\eps}\nabla w_\eps}_2\le CH\).
Also \(\norm{(w_\eps)_{j+1}}_2\le\norm{w_\eps}_2\le H\).
Equation~\eqref{eq:localized-energy} therefore bounds
\(\norm{\nabla(w_\eps)_j}_2\) by \(C_\lambda H\), independently
of \(\eps\) and \(j\).

We have \(w_\eps\to h\) in \(L^2\): use
Lemma~\ref{lem:smoothing} and the shrinking area of the region where
\(\chi_\eps\ne1\). For each fixed \(j\), choose a weakly convergent
subsequence of the globally bounded coefficient gradients. Testing in
regular charts identifies its limit as the distributional gradient
of \(h_j\). Weak lower semicontinuity proves
\eqref{eq:bounded-modes}.
\end{proof}

This estimate controls one coefficient at a time; it does not yet
give summability of their energies over \(j\). Its hypothesis is only
the local equation \(Ah=0\), so it will apply to weak limits without
requiring a separate flow-invariance argument. The next section uses
long orbit averages to obtain the additional identities needed for
the sum over all coefficients.

\section{Transverse rigidity of bounded eigenfunctions}\label{sec:rigidity}

The coefficient estimates of Section~\ref{sec:energy} do not yet give
square summability of the horizontal derivatives.  We obtain that summability
from an identity involving the two modes at the lower end of an angular
Fourier tail.
The long averages used below adapt the construction for invariant functions
in \cite[Sections~4--5]{Companion}; the phase in the average and the endpoint
terms retain the spectral parameter throughout the argument.

\begin{proof}[Proof of Proposition~\ref{prop:rigidity}]
Write $H=\norm{f}_\infty$ and $A=X-i\lambda$.  The eigenfunction equation
implies $Af=0$ locally in distributions.  Lemma~\ref{lem:bounded-modes}
therefore gives
\begin{equation}\label{eq:rigidity-mode-bound}
 \norm{\nabla f_j}_2\le C_\lambda H
 \qquad(j\in\Z).
\end{equation}
Here and below constants may depend on the fixed triangle, but not on an
angular index or on the approximation parameters.

\paragraph{The two boundary modes of a Fourier tail.}
Fix $m\in\Z$.  Define the $L^2$ field and the two coefficient derivatives
\[
 U=\sum_{j\ge m}P_jf,\qquad
 b=\pb f_m,\qquad a=\pa f_{m-1},
\]
and put
\[
 B=b(x)e^{i(m-1)\theta}-a(x)e^{im\theta}.
\]
The factors in this expression respect the coefficient transition rules,
so $B$ is a global $L^2(SM)$ field.  The mode equations give
\begin{equation}\label{eq:tail-equation}
 AU=B\qquad\text{locally in distributions on }SM.
\end{equation}
For completeness, for $N\ge m+2$ the finite tail satisfies
\[
 A\sum_{j=m}^N P_jf
 =B-(\pb f_{N+1})e^{iN\theta}
      +(\pa f_N)e^{i(N+1)\theta}.
\]
The last two terms are bounded in $L^2$ by
\eqref{eq:rigidity-mode-bound} and converge weakly to zero: their angular
frequencies tend to infinity.  Since the finite tails converge strongly
to $U$ in $L^2$, this proves \eqref{eq:tail-equation}.

We will prove the endpoint identity
\begin{equation}\label{eq:endpoint-identity}
 \ip{B}{(P_{m-1}+P_m)Yf}=0.
\end{equation}
At this stage $Yf$ denotes a distribution; its two displayed modes are
in $L^2$ by \eqref{eq:modes} and \eqref{eq:rigidity-mode-bound}.
The difficulty is that $Yf$ itself is not yet known to be an admissible
test in \eqref{eq:tail-equation}.  We construct compactly supported
approximations whose error under $A$ becomes negligible in that pairing.

\paragraph{Long averages away from the vertices.}
Let $K$ be the constant in Lemma~\ref{lem:smoothing} and fix $L>K+4$.
For $0<s<1$ and sufficiently small $\eps>0$, set $T=s/\eps$.
Let $G_\eps$ consist of the regular states $z$ whose entire trajectory
segment $\{\Psi_tz:0\le t\le T\}$ lies over $\{r>L\eps\}$.
Then
\begin{equation}\label{eq:bad-segments}
 \nu(SM\setminus G_\eps)\le C(s+\eps^2).
\end{equation}
Indeed, $r$ is $1$-Lipschitz along a unit-speed geodesic.  If a segment
meets $\{r\le L\eps\}$, a time mesh of spacing at most $\eps$ detects
a visit to $\{r\le(L+1)\eps\}$.  There are at most $T/\eps+2$ mesh
points.  Measure preservation and the cone-area bound
$\nu\{r\le(L+1)\eps\}\le C\eps^2$ prove
\eqref{eq:bad-segments}.  The discarded trajectories that hit a vertex
form a null set.

Define
\begin{equation}\label{eq:segment-average}
 q_\eps=f\ind_{G_\eps},\qquad
 g_\eps=\frac1T\int_0^T e^{i\lambda t}
                      q_\eps\circ\Psi_{-t}\,dt.
\end{equation}
The integral may be taken in $L^2$; bounded measurable representatives
and Fubini give the same field.  The phase is chosen so that replacing
$q_\eps$ by $f$ in the integrand gives $f$.  Consequently,
\begin{equation}\label{eq:average-bounds}
 \norm{g_\eps}_\infty\le H,\qquad
 \norm{g_\eps-f}_2\le CH\sqrt{s+\eps^2},\qquad
 \norm{Ag_\eps}_2\le\frac{2H}{T}.
\end{equation}
The first bound follows directly from \eqref{eq:segment-average}; the
second follows by integrating the $L^2$ distance
$\norm{q_\eps-f}_2\le H\nu(G_\eps^c)^{1/2}$.
For the third we use the distributional identity
\begin{equation}\label{eq:average-endpoints}
 Ag_\eps=\frac1T
       \bigl(q_\eps-e^{i\lambda T}q_\eps\circ\Psi_{-T}\bigr).
\end{equation}
Here no derivative of $\ind_{G_\eps}$ is taken.  To verify the identity,
test against $\varphi\in C_c^\infty(SM)$ and change variables by the
flow.  Along every regular orbit,
\[
 \frac{d}{dt}\left(e^{i\lambda t}
                  \overline{\varphi(\Psi_tz)}\right)
 =e^{i\lambda t}\overline{(A\varphi)(\Psi_tz)}.
\]
Integrating this formula in $t$, and using the formal skew adjointness
of $A$, proves \eqref{eq:average-endpoints}.  The test and its horizontal
derivative are bounded, which also justifies Fubini in this calculation.

Both $g_\eps$ and $Ag_\eps$ vanish over $\{r<L\eps\}$.  In fact, for
each $0\le t\le T$ the support of
$q_\eps\circ\Psi_{-t}$ is contained, up to null sets, in
$\Psi_tG_\eps\subset\{r>L\eps\}$; this also applies to both endpoint
terms in \eqref{eq:average-endpoints}.
Apply the spatial smoothing of Lemma~\ref{lem:smoothing} and set
\[
 w_\eps=S_\eps g_\eps.
\]
Before its prescribed zero extension this field vanishes on
$\{K\eps<r<(L-1)\eps\}$.  Thus the extension is smooth in position,
has compact spatial support in $M$, and creates no boundary derivative.
Local commutation with $A$ gives, globally as an $L^2$ identity,
$Aw_\eps=S_\eps Ag_\eps$.  The smoothing estimates imply
\begin{equation}\label{eq:regularized-bounds}
 \begin{aligned}
 \norm{w_\eps}_\infty&\le H,
 &\norm{Aw_\eps}_2&\le\frac{CH}{T},\\
 \norm{\nabla w_\eps}_2&\le\frac{CH}{\eps},
 &\norm{YAw_\eps}_2&\le\frac{CH}{T\eps}=\frac{CH}{s}.
 \end{aligned}
\end{equation}
All higher spatial derivatives also belong to $L^2$ for fixed $\eps$.
The relation $T\eps=s$ bounds the discarded measure by $C(s+\eps^2)$
while keeping $YAw_\eps$ bounded in $L^2$ at fixed $s$. Its norm need
not tend to zero: weak convergence suffices for pairing against the fixed
tail $U$. We first let $\eps$ tend to zero at fixed $s$ to obtain exact
bounded local solutions. Only then can Lemma~\ref{lem:bounded-modes}
restore coefficient estimates independent of $s$ and allow the limit
$s\to0$.

\paragraph{The first limit: $\eps\to0$ at fixed $s$.}
The last bound in \eqref{eq:regularized-bounds} makes $YAw_\eps$ bounded
in $L^2$ at fixed $s$.  Since $Aw_\eps\to0$ in $L^2$, for each
$\varphi\in C_c^\infty(SM)$ we have
\[
 \ip{YAw_\eps}{\varphi}
 =-\ip{Aw_\eps}{Y\varphi}\longrightarrow0.
\]
Density of these tests in $L^2(SM)$ therefore yields
\begin{equation}\label{eq:weak-error}
 YAw_\eps\rightharpoonup0\quad\text{in }L^2(SM).
\end{equation}
We may now test \eqref{eq:tail-equation} against $Yw_\eps$ and commute
the horizontal fields:
\begin{equation}\label{eq:tail-pairing}
 \ip{B}{Yw_\eps}
 =-\ip{U}{AYw_\eps}
 =-\ip{U}{YAw_\eps}\longrightarrow0.
\end{equation}
This use of a test with only $L^2$ angular regularity is legitimate.
Indeed, the finite angular Fourier sums of $Yw_\eps$ are smooth compactly
supported tests.  They converge to $Yw_\eps$, together with their first
horizontal derivatives, in $L^2$, because $w_\eps$ has square-integrable
spatial derivatives of every order.  Apply
\eqref{eq:tail-equation} first to these finite sums and then pass to the
limit.  The equality $AYw_\eps=YAw_\eps$ follows from $[X,Y]=0$.

To identify the left side of \eqref{eq:tail-pairing} after taking a
limit, we also need coefficient-gradient bounds.  Lemma~\ref{lem:energy},
with $E=M$, and \eqref{eq:regularized-bounds} give
\begin{equation}\label{eq:fixed-s-mode-bound}
 \norm{\nabla(w_\eps)_j}_2^2
 \le CH^2(s^{-1}+\lambda^2)
 \qquad(j\in\Z).
\end{equation}
Choose a sequence $\eps\downarrow0$ along which $w_\eps$ converges weakly
in $L^2$ to a field $h_s$.  Then
\begin{equation}\label{eq:first-limit}
 \norm{h_s}_\infty\le H,
 \qquad Ah_s=0\text{ locally in distributions},
 \qquad \norm{h_s-f}_2\le CH\sqrt{s}.
\end{equation}
For the first assertion use weak closedness in $L^2$ of the convex set
$\{u:|u|\le H\ \text{a.e.}\}$; the second follows from
$Aw_\eps\to0$ in $L^2$.  For the third, write
\[
 w_\eps-f=S_\eps(g_\eps-f)+(S_\eps f-f)
\]
and use \eqref{eq:average-bounds}, uniform $L^2$ boundedness of
$S_\eps$, and $S_\eps f\to f$ in $L^2$.

Only the coefficients with indices $m-2,m-1,m,m+1$ contribute to the
two modes $(P_{m-1}+P_m)Yw_\eps$.  By
\eqref{eq:fixed-s-mode-bound}, a further subsequence makes their
gradients converge weakly in the corresponding global $L^2$ coefficient
spaces.  Testing in regular coordinate charts identifies these limits
as the gradients of the respective coefficients of $h_s$.
Since $B$ has only modes $m-1$ and $m$, \eqref{eq:tail-pairing} gives
\begin{equation}\label{eq:first-limit-pairing}
 \ip{B}{(P_{m-1}+P_m)Yh_s}=0.
\end{equation}

\paragraph{The second limit: $s\to0$.}
The bound \eqref{eq:fixed-s-mode-bound} deteriorates as $s\to0$.
The first limit has replaced the approximate equation for $w_\eps$
by the exact local equation $Ah_s=0$, while preserving boundedness.
Lemma~\ref{lem:bounded-modes} now gives the stronger estimate
\[
 \norm{\nabla(h_s)_j}_2\le C_\lambda H
 \qquad(j\in\Z),
\]
independently of $s$.  No flow-invariance assertion for $h_s$ is needed
here: the lemma requires only its boundedness and its local
distributional equation.

Choose any sequence $s\downarrow0$.  By \eqref{eq:first-limit},
$h_s\to f$ strongly in $L^2$.  After passing to a subsequence, the
gradients of the same four coefficients converge weakly in $L^2$;
local distributional testing identifies their limits as those of $f$.
Passing to the limit in \eqref{eq:first-limit-pairing} proves
\eqref{eq:endpoint-identity}.  The successive limits have thus removed
the approximation error without losing the coefficient derivatives
needed by the endpoint pairing.

\paragraph{The sharp horizontal-energy bound.}
By $Af=0$ and the mode formulas \eqref{eq:modes},
\[
 (Yf)_{m-1}=-\lambda f_{m-1}-2ib,\qquad
 (Yf)_m=\lambda f_m+2ia.
\]
Substitute these expressions into \eqref{eq:endpoint-identity}.
With inner products linear in the first argument, the result is
\[
 2i\bigl(\norm{b}_2^2+\norm{a}_2^2\bigr)
 =\lambda\bigl(\ip{b}{f_{m-1}}+\ip{a}{f_m}\bigr).
\]
Taking absolute values and applying Cauchy--Schwarz gives
\begin{equation}\label{eq:endpoint-energy}
 \norm{\pb f_m}_2^2+\norm{\pa f_{m-1}}_2^2
 \le\frac{\lambda^2}{4}
       \bigl(\norm{f_{m-1}}_2^2+\norm{f_m}_2^2\bigr).
\end{equation}
This remains valid when the left side is zero; otherwise divide the
Cauchy--Schwarz inequality by its square root.
Summing \eqref{eq:endpoint-energy} over all $m\in\Z$ yields
\begin{equation}\label{eq:sharp-energy}
 \sum_{j\in\Z}\norm{\nabla f_j}_2^2
 =2\sum_{j\in\Z}
     \bigl(\norm{\pa f_j}_2^2+\norm{\pb f_j}_2^2\bigr)
 \le\lambda^2\norm{f}_2^2.
\end{equation}
In particular the horizontal derivatives of the angular partial sums
converge in $L^2$; their limits are the weak derivatives of $f$, by
local testing.  Parseval and \eqref{eq:gradient} therefore give
\[
 \norm{Xf}_2^2+\norm{Yf}_2^2
 =\norm{\nabla f}_2^2
 \le\lambda^2\norm{f}_2^2.
\]
But $Xf=i\lambda f$, so the first term already equals the right side.
It follows that $Yf=0$ and equality holds throughout.
\end{proof}

\section{Side reflections and weak mixing}\label{sec:conclusion}

The analytic argument is complete. We now apply Proposition~\ref{prop:rigidity}
to the lift of a billiard eigenfunction. Its local form will make the angle
condition decisive.

\begin{proposition}\label{prop:eigenfunctions}
Let \(Q\) be a nondegenerate Euclidean triangle with an angle irrational
relative to \(\pi\). If \(F\in L^2(Q\times S^1,\mu)\) satisfies
\[
 F\circ\Phi_t=e^{i\lambda t}F\quad\text{in }L^2
 \qquad(t\in\R)
\]
for some \(\lambda\in\R\), then \(F\) is constant almost everywhere.
If this constant is nonzero, then \(\lambda=0\).
\end{proposition}

\begin{proof}
First suppose \(F\) is bounded. Put the origin of the triangle coordinates
at a vertex whose angle \(\alpha\) satisfies \(\alpha/\pi\notin\mathbb Q\).
Let \(f\) be the lift of \(F\) to the double, with its orientation chosen
so that these coordinates are oriented on the first sheet.
Proposition~\ref{prop:rigidity} and the eigenfunction equation give
\[
 Xf=i\lambda f,\qquad Yf=0,
 \qquad\text{hence}\qquad \nabla f=i\lambda v f
\]
in every oriented flat chart, where \(v=(\cos\theta,\sin\theta)\).
Thus \(e^{-i\lambda x\cdot v}f(x,v)\) has zero spatial distributional
gradient. On a connected chart it is independent of position for almost
every direction. One can see this by taking its angular Fourier coefficients:
each has zero spatial distributional gradient and is therefore constant
on the connected base chart. In the first triangle interior this gives
\begin{equation}\label{eq:plane-wave}
 F(x,v)=c(v)e^{i\lambda x\cdot v}
\end{equation}
for a measurable bounded function \(c\) on \(S^1\).

Consider an open side. Reflection in its supporting line is the affine map
\(p\mapsto Rp+d\), where \(R\) is the linear orthogonal reflection and
\(d\in\R^2\). A crossing chart agrees with the first-sheet coordinates
and assigns coordinates \(x=Rp+d\) to the point \(p\) of the second sheet.
The chart direction \(v\) folds to the billiard direction \(Rv\) there.
Since \(R^T=R=R^{-1}\), Equation~\eqref{eq:plane-wave} gives
\[
 f(x,v)=F(p,Rv)=c(Rv)e^{i\lambda (x-d)\cdot v}
\]
on the second portion of the chart. The phase-removed function is independent
of position throughout the connected crossing chart, so its expressions on
the two open portions agree:
\begin{equation}\label{eq:side-relation}
 c(v)=c(Rv)e^{-i\lambda d\cdot v}
 \quad\text{for almost every }v\in S^1.
\end{equation}
This comparison uses open subsets of a regular chart and requires no
boundary trace for \(F\).

For the two sides through the origin, \(d=0\). Their reflections generate
a rotation through \(2\alpha\), under which \(c\) is invariant.
If \(\widehat c(k)\) is its \(k\)-th circle Fourier coefficient, this gives
\(\widehat c(k)=e^{2ik\alpha}\widehat c(k)\). Irrationality of
\(\alpha/\pi\) forces every coefficient except \(k=0\) to vanish.
Thus \(c\) is constant almost everywhere.

The third side has \(d\ne0\), since its supporting line does not pass
through the chosen vertex. If the constant \(c\) is nonzero,
\eqref{eq:side-relation} says \(e^{-i\lambda d\cdot v}=1\) almost
everywhere on the circle, and hence everywhere by continuity. The values
of \(d\cdot v\) fill the interval \([-|d|,|d|]\), so this is possible
only when \(\lambda=0\). Equation~\eqref{eq:plane-wave} now makes \(F\)
constant, proving the bounded case.

For an arbitrary \(L^2\) eigenfunction, use
\(\widetilde F=F/(1+|F|)\). The map \(z\mapsto z/(1+|z|)\) is bounded,
injective, and commutes with multiplication by unit complex scalars.
Consequently \(\widetilde F\) is a bounded eigenfunction with the same
frequency. The bounded case and injectivity give the assertion for \(F\).
\end{proof}

We recall the standard compact-operator proof of the spectral implication
needed to finish the argument; see, for example,
\cite[Theorem~6.1.4]{Peterson2015}. It applies
to the full continuous-time flow.

\begin{lemma}\label{lem:spectral}
Let \((\Omega,\mu)\) be a standard probability space and \(\Phi_t\) a
measure-preserving flow whose composition operators form a strongly
continuous unitary group on complex \(L^2(\mu)\). If the only flow
eigenfunctions are constants, then the product flow
\(\Phi_t\times\Phi_t\) is ergodic for \(\mu\otimes\mu\).
\end{lemma}

\begin{proof}
Write \(U_t h=h\circ\Phi_t\). The frequency-zero hypothesis first implies
ergodicity of \(\Phi_t\). If the product flow were not ergodic, there
would be a nonzero invariant kernel \(K\in L^2(\mu\otimes\mu)\) of mean
zero. Each one-variable marginal of \(K\) is invariant under the single
flow and thus constant. The zero total mean makes both marginals zero.

Define the complex-linear Hilbert--Schmidt operator
\[
 (Th)(x)=\int_\Omega K(x,y)h(y)\,d\mu(y).
\]
The zero marginals imply \(T1=0\) and that its range is contained in
\(L^2_0(\mu)=\{h:\int h\,d\mu=0\}\). Its restriction to this space
is nonzero, since a zero integral operator has zero \(L^2\) kernel.
Product invariance gives \(U_tT=TU_t\) by a change of variables in the
integral. Hence \(T^*T\), acting on \(L^2_0(\mu)\), is a nonzero
positive compact operator commuting with every \(U_t\).

A positive eigenvalue of \(T^*T\) has a nonzero finite-dimensional
eigenspace invariant under every \(U_t\). Commuting unitary operators
on that space have a common eigenvector \(h\ne0\). The associated
scalar character is continuous by strong continuity of \(U_t\), so
it has the form \(e^{i\lambda t}\) for some \(\lambda\in\R\).
This gives a flow eigenfunction in \(L^2_0(\mu)\), a contradiction.
\end{proof}

\begin{proof}[Proof of Theorem~\ref{thm:main}]
The billiard composition group is strongly continuous by the construction
in Section~\ref{sec:geometry}. Proposition~\ref{prop:eigenfunctions}
eliminates every nonconstant eigenfunction, so Lemma~\ref{lem:spectral}
proves ergodicity of the product flow.
\end{proof}

\bibliographystyle{plain}
\bibliography{references}
\end{document}
